\documentclass[12pt]{article}
\usepackage[utf8]{inputenc}
\usepackage[T1]{fontenc}
\usepackage[margin=1in]{geometry}
\usepackage{setspace}
\usepackage{xcolor}
\usepackage[colorlinks=true, urlcolor=blue!60!black, linkcolor=blue!60!black]{hyperref}
\setlength{\parindent}{2em}
\setlength{\parskip}{1.25em}
\renewcommand{\baselinestretch}{1.0}
\begin{document}
\begin{center}
{\LARGE \textit{Plessy v.\ Ferguson}, 163 U.S.\ 537 (1896)\par}
\vspace{0.8em}
{\large Case Brief\par}
\end{center}
\vspace{1em}

\section{Memory Jogger}

By a 7--1 vote over Justice Harlan's lone ``color-blind Constitution'' dissent, the Supreme Court upheld Louisiana's Separate Car Act under the ``separate but equal'' doctrine, ruling that state-mandated racial segregation in intrastate rail travel was a reasonable exercise of the police power that violated neither the Thirteenth nor the Fourteenth Amendment.

\section{Detailed Case Facts}

On July 10, 1890, the Louisiana General Assembly enacted the Separate Car Act (Acts 1890, No.\ 111, p.\ 152), requiring all railway companies carrying passengers within Louisiana (excluding street railroads) to provide ``equal but separate accommodations for the white, and colored races,'' either by providing two or more passenger coaches or by dividing coaches with a partition (\href{https://www.law.cornell.edu/supremecourt/text/163/537}{Cornell LII}; \href{https://www.courtlistener.com/opinion/94508/plessy-v-ferguson/}{CourtListener}). Section 2 required train officers to assign passengers to the coach designated for their race and made any passenger who insisted on occupying a coach assigned to another race liable to a \$25 fine or up to 20 days' imprisonment in the parish prison. Section 3 exempted only ``nurses attending children of the other race.''

To challenge the statute, a New Orleans civil-rights organization, the \textit{Comit\'e des Citoyens} (Citizens' Committee), organized a deliberate test case with the quiet cooperation of the East Louisiana Railway, which opposed the expense of maintaining separate cars (\href{https://supreme.justia.com/cases/federal/us/163/537/}{Justia}). The Committee recruited Homer Adolph Plessy, a U.S.\ citizen and Louisiana resident of mixed descent---seven-eighths Caucasian and one-eighth African ancestry, in whom ``the mixture of colored blood was not discernible''---and engaged a private detective to ensure Plessy would be arrested specifically under the Separate Car Act rather than for disturbing the peace (\href{https://www.law.cornell.edu/supremecourt/text/163/537}{Cornell LII}; \href{https://supreme.justia.com/cases/federal/us/163/537/}{Justia}).

On June 7, 1892, Plessy purchased a first-class ticket on the East Louisiana Railway for an intrastate trip from New Orleans to Covington, Louisiana, boarded the train, and took a vacant seat in the whites-only coach. When ordered by the conductor to vacate the seat and move to the coach assigned to Black passengers, Plessy refused. With the assistance of a police officer, he was forcibly ejected from the train, arrested, and jailed in the Orleans Parish prison to answer a criminal charge of violating the Separate Car Act (\href{https://www.law.cornell.edu/supremecourt/text/163/537}{Cornell LII}).

\section{Procedural History}

\begin{itemize}
\item \textbf{Criminal District Court for the Parish of Orleans (1892):} Following preliminary examination before the city recorder, a criminal information was filed against Plessy before Judge John Howard Ferguson. Plessy interposed a plea asserting that the Separate Car Act was null and void under the Thirteenth and Fourteenth Amendments to the U.S.\ Constitution. The district attorney filed a demurrer, which Judge Ferguson sustained---overruling Plessy's plea and ordering him to plead over to the facts, distinguishing intrastate travel from interstate travel (\href{https://www.law.cornell.edu/supremecourt/text/163/537}{Cornell LII}).
\item \textbf{Supreme Court of Louisiana (1892):} Because no direct appeal lay from a sentence under the statute, Plessy petitioned the Supreme Court of Louisiana for writs of prohibition and certiorari to enjoin Judge Ferguson from proceeding with the criminal prosecution. The state supreme court held the Separate Car Act constitutional as applied to intrastate rail travel and denied relief (\textit{Ex parte Plessy}, 45 La.\ Ann.\ 80, 11 So.\ 948 (1892)) (\href{https://www.law.cornell.edu/supremecourt/text/163/537}{Cornell LII}).
\item \textbf{U.S.\ Supreme Court (1896):} The Chief Justice of the Supreme Court of Louisiana allowed a writ of error to the U.S.\ Supreme Court. Represented by Albion W.\ Tourg\'ee and Samuel F.\ Phillips, Plessy argued the case on April 13, 1896; on May 18, 1896, the Supreme Court affirmed the judgment of the Supreme Court of Louisiana by a 7--1 vote (\href{https://www.oyez.org/cases/1850-1900/163us537}{Oyez}; \href{https://www.archives.gov/milestone-documents/plessy-v-ferguson}{National Archives}).
\item \textbf{Subsequent History:} The ``separate but equal'' doctrine legitimized Jim Crow segregation across the American South for nearly six decades until the Supreme Court rejected it in public education in \textit{Brown v.\ Board of Education}, 347 U.S.\ 483 (1954), and subsequently dismantled segregation in public transportation (e.g., \textit{Gayle v.\ Browder}, 352 U.S.\ 903 (1956)) and other public facilities (\href{https://supreme.justia.com/cases/federal/us/163/537/}{Justia}).
\end{itemize}

\section{Judicial Votes}

The Court decided \textit{Plessy v.\ Ferguson} by a \textbf{7--1} vote, with one Justice not participating (\href{https://www.oyez.org/cases/1850-1900/163us537}{Oyez}; \href{https://www.law.cornell.edu/supremecourt/text/163/537}{Cornell LII}):
\begin{itemize}
\item \textbf{Majority (7):} Justice Henry Billings Brown (author of the opinion of the Court), joined by Chief Justice Melville W.\ Fuller and Justices Stephen J.\ Field, Horace Gray, George Shiras Jr., Edward D.\ White, and Rufus W.\ Peckham.
\item \textbf{Dissent (1):} Justice John Marshall Harlan.
\item \textbf{Did Not Participate (1):} Justice David J.\ Brewer did not hear oral argument or participate in the decision because he had returned home to Kansas following the sudden death of his daughter (\href{https://www.law.cornell.edu/supremecourt/text/163/537}{Cornell LII}; \href{https://en.wikipedia.org/wiki/David_J._Brewer}{Wikipedia -- David J.\ Brewer}).
\end{itemize}

\section{Holding}

\textbf{Judgment:} Affirmed. The Supreme Court upheld the judgment of the Supreme Court of Louisiana denying Plessy's petition for writs of prohibition and certiorari, sustaining the constitutionality of Louisiana's Separate Car Act as applied to intrastate rail transportation (\href{https://www.law.cornell.edu/supremecourt/text/163/537}{Cornell LII}).

\textbf{Rule Announced:} State laws mandating ``equal but separate'' accommodations for white and Black passengers in intrastate public conveyances do not violate the Thirteenth Amendment (which forbids only actual slavery or involuntary servitude) or the Fourteenth Amendment's Equal Protection Clause (which secures political and civil equality before the law, not ``social'' equality), so long as the racial separation is a ``reasonable'' exercise of the State's police power enacted in good faith with reference to ``the established usages, customs, and traditions of the people'' (\href{https://www.law.cornell.edu/supremecourt/text/163/537}{Cornell LII}).

\textbf{Holding vs.\ Dicta:} The Court's holding was limited to upholding Louisiana's requirement of separate-but-equal accommodations on intrastate railways against Thirteenth and Fourteenth Amendment challenges. Several influential passages were technically \textit{dicta}, including: (1) the Court's endorsement of racially segregated public schools and anti-miscegenation laws as valid exercises of state police power; (2) its suggestion that the statutory clause immunizing railroads from damages for misclassifying a white passenger might be unconstitutional (a question not presented on the record); (3) its statement that determining the proportion of ancestry that defines race is a matter of state law; and (4) its broad sociological assertions that ``[l]egislation is powerless to eradicate racial instincts'' and that any ``badge of inferiority'' existed ``solely because the colored race chooses to put that construction upon it'' (\href{https://www.law.cornell.edu/supremecourt/text/163/537}{Cornell LII}).

\section{Analysis of Opinions}

\subsection{Majority Opinion (Justice Brown)}

\textbf{Thirteenth Amendment:} Justice Brown dismissed the Thirteenth Amendment claim as ``too clear for argument.'' Citing the \textit{Slaughter-House Cases}, 83 U.S.\ (16 Wall.) 36 (1873), and the \textit{Civil Rights Cases}, 109 U.S.\ 3 (1883), he defined slavery narrowly as legal bondage or involuntary servitude. A statute drawing ``merely a legal distinction between the white and colored races,'' he reasoned, neither destroyed legal equality nor re-established involuntary servitude (\href{https://www.law.cornell.edu/supremecourt/text/163/537}{Cornell LII}).

\textbf{Fourteenth Amendment---Political vs.\ Social Equality:} Turning to the Equal Protection Clause, Brown acknowledged that the Amendment aimed to enforce ``the absolute equality of the two races before the law,'' but asserted that ``in the nature of things, it could not have been intended to abolish distinctions based upon color, or to enforce social, as distinguished from political, equality, or a commingling of the two races upon terms unsatisfactory to either.'' As proof that racial separation was widely accepted under the police power, he pointed to racially segregated public schools---including Chief Justice Shaw's pre-Fourteenth Amendment decision in \textit{Roberts v.\ City of Boston}, 59 Mass.\ (5 Cush.) 198 (1850), and Congress's segregated schools in the District of Columbia---as well as state anti-miscegenation statutes. He distinguished \textit{Strauder v.\ West Virginia}, 100 U.S.\ 303 (1880) (striking down exclusion of Black citizens from juries), as involving political and civil rights rather than social accommodations, and noted that because the East Louisiana Railway operated wholly within Louisiana, no Dormant Commerce Clause issue arose under \textit{Louisville, New Orleans \& Texas Railway Co.\ v.\ Mississippi}, 133 U.S.\ 587 (1890) (\href{https://www.law.cornell.edu/supremecourt/text/163/537}{Cornell LII}).

\textbf{Property in Whiteness and ``Reasonableness'' Review:} Addressing Plessy's due process argument that the reputation of being white was a valuable ``property'' right, Brown replied that if Plessy were legally white he could sue the railroad for damages, whereas if he were legally ``colored'' under Louisiana law, he was not lawfully entitled to the reputation of being white. Responding to Tourg\'ee's slippery-slope hypotheticals (separate cars by hair color or separate sides of the street by race), Brown stated that police-power regulations must be ``reasonable'' and enacted in good faith rather than for ``the annoyance or oppression of a particular class'' (citing \textit{Yick Wo v.\ Hopkins}, 118 U.S.\ 356 (1886)). Yet he gave legislatures ``a large discretion'' to act in accordance with ``the established usages, customs, and traditions of the people'' to preserve ``public peace and good order.'' Finally, Brown rejected the claim that mandated separation stamped Black citizens with a badge of inferiority, blaming any such perception ``solely'' on the construction Black citizens chose to place on the law, and declaring that if one race were socially inferior to another, ``the constitution of the United States cannot put them upon the same plane'' (\href{https://www.law.cornell.edu/supremecourt/text/163/537}{Cornell LII}).

\subsection{Dissenting Opinion (Justice Harlan)}

Justice John Marshall Harlan delivered a solitary, prophetic dissent attacking both the legal reasoning and the social dishonesty of the majority opinion. First, he emphasized that a railroad is a ``public highway'' performing public functions, meaning that state-mandated exclusion on trains directly infringed citizens' civil rights and personal liberty of locomotion, not merely ``social'' intercourse (\href{https://www.law.cornell.edu/supremecourt/text/163/537}{Cornell LII}).

Second, Harlan rejected the majority's formalistic claim that the statute treated both races symmetrically. ``Every one knows,'' he wrote, that the Act's real purpose was ``not so much to exclude white persons from railroad cars occupied by blacks, as to exclude colored people from coaches occupied by or assigned to white persons.'' Such state-enforced separation imposed a ``badge of servitude'' forbidden by the Thirteenth Amendment and violated the Fourteenth Amendment's guarantee of equal citizenship. Warning that the majority's ``reasonableness'' standard would permit states to segregate streets, courtrooms, legislative galleries, or even jury boxes by race or religion, Harlan declared that ``[t]here is no caste here. Our constitution is color-blind, and neither knows nor tolerates classes among citizens.'' He predicted that the judgment would prove ``quite as pernicious as the decision made by this tribunal in the \textit{Dred Scott Case}'' and would plant ``the seeds of race hate'' under sanction of law (\href{https://www.law.cornell.edu/supremecourt/text/163/537}{Cornell LII}).

Notably, however, Harlan's dissent also reflected late-nineteenth-century racial hierarchies: immediately before proclaiming a color-blind Constitution, he affirmed that ``[t]he white race deems itself to be the dominant race in this country'' in ``prestige, in achievements, in education, in wealth, and in power,'' and later contrasted Black citizens' exclusion with the fact that a ``Chinaman''---whom Harlan described as belonging to ``a race so different from our own that we do not permit those belonging to it to become citizens of the United States''---could legally ride in the white coach (\href{https://www.law.cornell.edu/supremecourt/text/163/537}{Cornell LII}).

\section{Examples: Future Applications}

\begin{enumerate}
\item \textbf{Same-Side (Upheld Under \textit{Plessy}'s Framework):} A state enacts a statute requiring all intrastate public buses to maintain separate front and rear seating sections for white and Black riders, nominally providing identical seat construction and fares for both groups. \textit{Why it falls here:} Under \textit{Plessy}, intrastate public transportation falls within the State's police power to regulate ``social'' contact in accordance with local customs, so long as accommodations are formally equal.
\item \textbf{Same-Side (Upheld Under \textit{Plessy}'s Framework):} A state legislature establishes separate public parks and libraries for white and Black residents, asserting that separation preserves ``public peace and good order'' while budgeting comparable facilities for each race. \textit{Why it falls here:} Because \textit{Plessy} extended its police-power rationale beyond trains to public schools and other sites of public commingling, state-mandated segregation in public accommodations and recreation was routinely sustained under \textit{Plessy}'s deferential reasonableness test prior to 1954.
\item \textbf{Opposite-Side (Invalidated Even Under \textit{Plessy}):} A state statute bars Black citizens from serving on grand or petit juries, or prohibits Black citizens from owning or occupying real property in designated neighborhoods. \textit{Why it falls here:} \textit{Plessy} itself reaffirmed \textit{Strauder v.\ West Virginia}, drawing a line between ``social'' accommodations and core ``political and civil rights'' (such as jury service, property ownership, and contract rights, later enforced in \textit{Buchanan v.\ Warley}, 245 U.S.\ 60 (1917)).
\item \textbf{Opposite-Side (Invalidated Even Under \textit{Plessy}):} A state funds an accredited public law school exclusively for white residents while offering no in-state legal education for Black residents (or offering a patently inferior facility lacking comparable faculty, library resources, and prestige). \textit{Why it falls here:} Even taking \textit{Plessy}'s ``separate \textit{but equal}'' formula at face value, the failure to provide substantially equal facilities violates the Equal Protection Clause---the very doctrine the Supreme Court later used in \textit{Missouri ex rel.\ Gaines v.\ Canada}, 305 U.S.\ 337 (1938), and \textit{Sweatt v.\ Painter}, 339 U.S.\ 629 (1950), to erode segregation from within.
\item \textbf{Fence-Sitter (Close / Contested Case):} A state statute requires separate dining-car curtains on interstate passenger trains passing through the state, affecting passengers traveling across state lines as well as intrastate riders. \textit{Why it is a fence-sitter:} While \textit{Plessy} upheld racial separation as consistent with the Fourteenth Amendment for \textit{intrastate} travel, Justice Brown expressly relied on \textit{Hall v.\ DeCuir}, 95 U.S.\ 485 (1878), and \textit{Louisville, N.O.\ \& T.\ Ry.\ Co.\ v.\ Mississippi} to limit the ruling to internal state commerce; applying a state segregation mandate to interstate passengers directly triggers Dormant Commerce Clause invalidation (as the Court later held in \textit{Morgan v.\ Virginia}, 328 U.S.\ 373 (1946)), creating a sharp tension between \textit{Plessy}'s Equal Protection deference and federal uniformity over interstate commerce.
\end{enumerate}

\section{Critique}

\textbf{Anti-Caste and Doctrinal Critiques:} \textit{Plessy v.\ Ferguson} is universally ranked alongside \textit{Dred Scott} as part of the constitutional ``anti-canon.'' In his classic defense of \textit{Brown}, Charles L.\ Black, Jr.\ exposed the willful blindness of Justice Brown's ``badge of inferiority'' passage: everyone in the Jim Crow South knew as a matter of common social fact that segregation was a comprehensive system designed to enforce white supremacy and subordinate Black citizens, not a symmetrical accommodation (\href{https://scholarship.law.columbia.edu/cgi/viewcontent.cgi?article=3132&context=faculty_scholarship}{Columbia Law Scholarship -- Reading Charles Black}). Scholars also highlight the stark inconsistency between \textit{Plessy}'s extreme deference to state police-power restrictions on Black citizens' freedom of movement and the Fuller and \textit{Lochner}-era Court's simultaneous aggressive invalidation of economic regulations in the name of liberty of contract (\href{https://scholarworks.indianapolis.iu.edu/bitstream/1805/5526/1/hacker-2005-the-neutrality.pdf}{Hacker, The Neutrality Principle}).

\textbf{Originalist and Historical Critiques:} Originalist scholars have intensely debated whether \textit{Plessy} was faithful to the 1868 original understanding of the Fourteenth Amendment. Michael W.\ McConnell famously argued that the Reconstruction Congress's debates over the Civil Rights Act of 1875 demonstrate that a majority of the Fourteenth Amendment's framers understood the Amendment to forbid \textit{de jure} racial segregation in public accommodations and schools, making \textit{Plessy} wrong on originalist grounds as well (\href{https://digital.sandiego.edu/cgi/viewcontent.cgi?article=3165&context=sdlr}{San Diego Law Review -- citing McConnell}). Conversely, historians and scholars such as Michael J.\ Klarman and Earl M.\ Maltz have contended that the 1868 ratifying public generally accepted the tripartite distinction among civil, political, and social rights invoked by Justice Brown, suggesting that \textit{Plessy} reflected mainstream nineteenth-century racial attitudes and that rejecting \textit{Plessy} requires reading the Equal Protection Clause at a higher level of constitutional principle rather than narrow original expected applications (\href{https://digital.sandiego.edu/cgi/viewcontent.cgi?article=3165&context=sdlr}{San Diego Law Review -- citing Klarman}; \href{https://journals.library.wustl.edu/jurisprudence/article/2090/galley/18923/view/}{Washington University Jurisprudence Review}).

\textbf{Modern Debates Over Harlan's Dissent:} Contemporary scholars also contest the legacy of Justice Harlan's dissent. Conservative jurists invoke Harlan's ``color-blind Constitution'' aphorism to condemn all race-conscious government action, including affirmative action, whereas progressive scholars such as Goodwin Liu argue that Harlan's dissent was fundamentally an \textit{anti-caste} opinion directed against racial subordination rather than a formal ban on remedial race-consciousness (\href{https://journals.law.harvard.edu/lpr/wp-content/uploads/sites/89/2013/05/2.1_3_Liu.pdf}{Liu, Harvard Law \& Policy Review}). Meanwhile, Gabriel J.\ Chin's landmark study, \textit{The Plessy Myth}, cautions against hagiography by highlighting Harlan's explicit anti-Chinese rhetoric in \textit{Plessy} and his mixed record in Chinese-exclusion cases (\href{https://digitalcommons.wne.edu/cgi/viewcontent.cgi?article=1725&context=lawreview}{Western New England Law Review -- discussing Chin}).

\textbf{Competing Underlying Values:} \textit{Plessy} pitted the Reconstruction Amendments' transformative promise of equal national citizenship and individual liberty against federalism, deferential state police power, and entrenched local racial customs---demonstrating how formalistic legal categories (``social'' vs.\ ``political'' rights, ``equal'' but separate) can mask systemic subordination.

\section{Key Quotations}

\begin{quote}
``The object of the amendment was undoubtedly to enforce the absolute equality of the two races before the law, but, in the nature of things, it could not have been intended to abolish distinctions based upon color, or to enforce social, as distinguished from political, equality, or a commingling of the two races upon terms unsatisfactory to either.''
\end{quote}
\noindent --- Justice Brown (Opinion of the Court), 163 U.S.\ at 544 (\href{https://www.law.cornell.edu/supremecourt/text/163/537}{Cornell LII}; \href{https://www.courtlistener.com/opinion/94508/plessy-v-ferguson/}{CourtListener}).

\begin{quote}
``We consider the underlying fallacy of the plaintiff's argument to consist in the assumption that the enforced separation of the two races stamps the colored race with a badge of inferiority. If this be so, it is not by reason of anything found in the act, but solely because the colored race chooses to put that construction upon it.''
\end{quote}
\noindent --- Justice Brown (Opinion of the Court), 163 U.S.\ at 551 (\href{https://www.law.cornell.edu/supremecourt/text/163/537}{Cornell LII}).

\begin{quote}
``Legislation is powerless to eradicate racial instincts, or to abolish distinctions based upon physical differences, and the attempt to do so can only result in accentuating the difficulties of the present situation.''
\end{quote}
\noindent --- Justice Brown (Opinion of the Court), 163 U.S.\ at 551 (\href{https://www.law.cornell.edu/supremecourt/text/163/537}{Cornell LII}).

\begin{quote}
``But in view of the constitution, in the eye of the law, there is in this country no superior, dominant, ruling class of citizens. There is no caste here. Our constitution is color-blind, and neither knows nor tolerates classes among citizens. In respect of civil rights, all citizens are equal before the law.''
\end{quote}
\noindent --- Justice Harlan (dissenting), 163 U.S.\ at 559 (\href{https://www.law.cornell.edu/supremecourt/text/163/537}{Cornell LII}).

\begin{quote}
``In my opinion, the judgment this day rendered will, in time, prove to be quite as pernicious as the decision made by this tribunal in the Dred Scott Case.''
\end{quote}
\noindent --- Justice Harlan (dissenting), 163 U.S.\ at 559 (\href{https://www.law.cornell.edu/supremecourt/text/163/537}{Cornell LII}).

\end{document}
